\item \points{1b} \textbf{Adding HER to Bit Flipping}

With HER, the model is trained on the actual (state, goal, reward) tuples along with (state,
goal, reward) tuples where the goal has been relabeled after the fact. The goals are relabeled to the state that was \textit{actually reached} and the rewards should be relabeled correspondingly. In other words, we pretend that the state we reached was our goal all along. HER
gives us more examples of actions that lead to positive rewards, effectively doubling our
experience (since we use the original episode as well as the relabeled one). The reward
function for relabeled goals is the same as the environment reward function; for the bit
flipping environment, the reward is $-1$ if the state and goal vector do not match and 0 if
they do match.

There are three different variations of HER: \texttt{final}, \texttt{random}, and \texttt{future}. Suppose the data collected in an episode consists of the following (state, goal, reward) tuples: $\{(s_t, g, r_t)\}_{t=0}^T$, where $t$ indicates each time step in the episode. Given a (state, goal, reward) tuple $(s_i, g_i, r_i)$ each variation of HER relabels the goal $g_i$ differently:

\begin{itemize}
    \item \texttt{final}: The final state of the episode is used as the goal.
    
    \item \texttt{random}: A random state in the episode is used as the goal.
    
    \item \texttt{future}: A random future state of the episode is used as the goal. Specifically, if you want to relabel the goal in the tuple $(s_i, g_i, r_i)$, only states $s_j$ with $j > i$ can be used.
\end{itemize}

Implementation Notes:

\begin{itemize}
    \item Always use \texttt{copy()} when assigning an existing numpy array to a new variable. NumPy does not create a new copy by default.
    
    \item When choosing a new goal for relabelling, choose the state corresponding to \texttt{next\_state} from the experience.
\end{itemize}

More details of these three implementations are in the comments of \texttt{trainer.py}. Implement the three variations of HER in the function \texttt{update\_replay\_buffer} in \texttt{trainer.py}. You \textbf{do not} need to submit a plot for this.


